%% CNSNS submission — Communications in Nonlinear Science and Numerical Simulation
%% Special Issue: "Artificial Intelligence Methods for Inverse Problems"
%% Document class: elsarticle (Elsevier); compiled with pdflatex + bibtex.
\documentclass[preprint,12pt]{elsarticle}

\usepackage{amsmath,amssymb}
\usepackage{graphicx}
\graphicspath{{../figures/}}
\usepackage{booktabs}
\usepackage{array}
\usepackage{tabularx}
\usepackage{multirow}
\usepackage{algorithm}
\usepackage{algpseudocode}
\usepackage{siunitx}
\usepackage{xcolor}
\usepackage[hidelinks]{hyperref}

\newcommand{\bp}{\ensuremath{\,\mathrm{bp}}}
\newcommand{\todo}[1]{\textcolor{red}{[#1]}}

%% Journal information
\journal{Communications in Nonlinear Science and Numerical Simulation}

%% ---------------------------------------------------------------
%% revision tables are included only when make_tables.py has written them
\newif\ifrevtable
%% keep floats near their text instead of deferring them to the end
\renewcommand{\topfraction}{0.9}
\renewcommand{\bottomfraction}{0.6}
\renewcommand{\textfraction}{0.05}
\renewcommand{\floatpagefraction}{0.75}
\setcounter{topnumber}{3}
\setcounter{totalnumber}{5}
%% LaTeX's \input wrapper breaks booktabs rules when used inside a tabular;
%% the primitive does not.
\makeatletter
\newcommand{\revrows}[1]{\@@input tables/#1.tex }
\makeatother

\begin{document}

\begin{frontmatter}

\title{A Physics-Informed Neural Network Forward Solver for the Rough-Volatility
Calibration Inverse Problem: Dimension Scaling on the Markovian-Lifted PDE and a
Fast Warm-Start Hybrid}

\author[cu]{Aviral Sharma\corref{cor}}
\ead{aviral98765@gmail.com}
\author[sgt]{Kamal Preet Singh}
\author[ind]{Vipul Sharma}

\cortext[cor]{Corresponding author}
\affiliation[cu]{organization={Chandigarh University},
              city={Chandigarh}, country={India}}
\affiliation[sgt]{organization={SGT University},
              city={Gurugram}, country={India}}
\affiliation[ind]{organization={Independent Researcher}, country={India}}

%% ----- Highlights (3-5 bullets, max 85 characters each) --------
\begin{highlights}
\item PINN forward solver for the Markovian-lifted rough-Heston pricing PDE
\item Solver error shows no detectable growth in lift dimension; grid cost explodes
\item Explicit decomposition separates network solve error from lift model error
\item Parametric PINN warm-start cuts expensive forward-model calls about 9-fold
\item Accuracy limits mapped: error grows in rough, high vol-of-vol corners
\end{highlights}

%% ----- Abstract ------------------------------------------------
\begin{abstract}
Calibrating a rough-volatility model to an option surface is an inverse problem
whose fractional, non-Markovian variance admits no finite-dimensional forward
partial differential equation (PDE). A Markovian lift approximates the kernel by
$n$ exponentials and recovers an $(n{+}1)$-dimensional forward PDE, which we
solve with a physics-informed neural network (PINN) trained on its residual.
Measured against an independent lifted Monte Carlo at the same $n$, and separated
from the lift's own approximation error, the network's solve-error shows no
statistically detectable growth as $n$ rises from 4 to 32 (log-log slope
$0.09\pm0.08$, eight seeds), whereas a naive full-tensor finite-difference
grid---an illustrative bound, not a tuned solver---is infeasible by $n{=}8$. The
solve/lift decomposition shows whether to add factors or improve the network. A
parametric extension is too inaccurate to calibrate alone, but as a warm start
for a short exact-model refinement it recovers full calibration quality at $8.8$
times fewer expensive evaluations. On live Bitcoin options from Deribit, the
surface pins the correlation ($\rho\approx-0.43$) but identifies the Hurst
exponent only to an interval, $H\in[0.04,0.30]$, well inside the rough regime; a
realised-volatility estimate using no option data falls inside it. The fit is
$113$ vol\bp{} vega-weighted and $416$ unweighted, the latter dominated by
options under two weeks; beyond two months it is $52$ vol\bp{}.
\end{abstract}

\begin{keyword}
inverse problems \sep physics-informed neural networks \sep rough volatility
\sep Markovian lift \sep model calibration \sep curse of dimensionality \sep
scientific machine learning
\end{keyword}

\end{frontmatter}

% ================================================================
\section{Introduction}
\label{sec:intro}

Model calibration is the canonical inverse problem of mathematical finance: given
a finite set of noisy, arbitrage-constrained option prices, recover the
parameters of a forward model that reproduces them. Rough-volatility models make
this inverse problem unusually severe, because their \emph{forward} problem is
already pathological. Motivated by the empirical roughness of realised variance
\citep{gatheral2018volatility,bennedsen2022decoupling} and the steep
short-maturity skew that classical stochastic-volatility models such as
Heston's~\citep{heston1993closed} cannot reproduce, rough-volatility models drive the
variance with a fractional noise whose Hurst exponent $H<1/2$. That single
fact---whose lineage runs back to long-memory volatility
\citep{comte1998long}---makes the variance process neither Markovian nor a
semimartingale, so there is no finite-dimensional forward PDE to solve, let alone
invert \citep{gatheral2018volatility,eleuch2019characteristic}.

The standard engineering escape is the \emph{Markovian lift}: approximate the
fractional kernel by a sum of $n$ exponentials and recover a genuine
$(n{+}1)$-dimensional forward PDE \citep{abijaber2019multifactor,abijaber2019lifting}.
This turns an ill-posed forward problem into a well-posed but high-dimensional
one, and pushes the difficulty of the inverse problem onto the cost of repeatedly
solving that high-dimensional forward PDE inside a calibration loop---once per
candidate parameter vector, at every iteration of the outer optimiser. Any AI
method proposed as a forward operator for this inverse problem must therefore be
judged on two separate questions: how accurately does it solve the forward PDE it
is handed, and how does that accuracy behave as the lifted dimension $n$---the
knob that trades model fidelity against forward-problem cost---is turned up. Both
questions are usually left unanswered when a neural network is dropped into a
calibration pipeline as a black-box pricer.

We answer them directly for a physics-informed neural network (PINN) forward
operator. Rather than laying down a grid, a PINN represents the solution as a
neural network and trains it to satisfy the PDE at randomly sampled collocation
points \citep{raissi2019physics}, sidestepping the exponential blow-up that
kills finite differences in high dimensions, in the spirit of the deep Galerkin
and deep BSDE solvers \citep{sirignano2018dgm,han2018solving}. Because it is
differentiable in its inputs, the same trained network doubles as a
differentiable forward operator for gradient-based calibration once the
parameters are added as additional inputs. Our concrete contributions are as
follows.

\begin{enumerate}

\item \textbf{A convergence and stability study of the AI forward operator.} We
show that the PINN's own solve-error---the part of the total error under the
network's control---shows no statistically detectable growth as the lifted
dimension climbs from 4 to 32 over eight seeds per dimension, and costs the same
per training iteration at every $n$, while a naive full-tensor
finite-difference forward solver ($50^{n+1}$ points) is already infeasible by
$n{=}8$. We identify the four training design choices without
which this convergence behaviour does not hold (Section~\ref{sec:arch}).

\item \textbf{An error decomposition that isolates what limits the inverse
problem.} We split total pricing error into the network's own solve-error
(measured against a lifted Monte Carlo at the same $n$) and the forward model's
lift-error (measured against the exact, non-lifted Fourier forward map). Their
crossover gives a concrete, maturity-dependent rule for how large $n$---and
hence how expensive the forward operator---needs to be before the network, not
the model, is the accuracy bottleneck.

\item \textbf{A differentiable forward operator embedded in the calibration
inverse problem.} A parametric extension of the network, additionally
conditioned on the five rough-Heston parameters, is used both as a fast
stand-alone forward map and as a warm start for a short local refinement against
the true forward model, and we report where its accuracy generalises and where
it does not.

\item \textbf{A real-market solution of the inverse problem.} Calibrated to live
Deribit Bitcoin option data, the surface identifies the correlation sharply and
the Hurst exponent only to an interval that lies well inside the rough regime.
We report that non-identification, its mechanism and its consequences for risk
functionals, rather than a point estimate the data do not support.

\item \textbf{Fast recalibration, quantified against the true forward model.}
The parametric warm start plus a short exact-model polish recovers from-scratch
calibration quality at $8.8$ times fewer expensive forward-model evaluations,
whichever parameter box the network was trained on; we state plainly where the standalone AI forward operator is not
accurate enough to trust alone.

\end{enumerate}

We are explicit about what is \emph{not} new.
Papapantoleon and Rou~\citep{papapantoleon2025timestepping} solved the lifted rough-Heston forward PDE
with a mesh-free network before us. Our contribution is a residual PINN forward
operator characterised specifically as a component of the calibration inverse
problem: a quantitative convergence-versus-dimension study, an error
decomposition that separates network error from model error, a real-market
inverse-problem solution, and the warm-start hybrid.

% ================================================================
\section{Related work}
\label{sec:related}

\paragraph{Closest prior work}
Papapantoleon and Rou~\citep{papapantoleon2025timestepping} price European options in rough-diffusion
models by recasting the forward pricing PDE as an energy-minimisation problem and
approximating it in a time-stepping deep gradient flow, demonstrated on the
lifted Heston model. Their method is the closest to ours in target and spirit; it
differs in mechanism (time-stepping variational scheme versus residual PINN) and
does not study forward-operator dimension scaling, error decomposition, or the
calibration inverse problem downstream.

\paragraph{Neural forward operators for rough volatility}
Jacquier and \v{Z}uri\v{c}~\citep{jacquier2023random} cast the path-dependent pricing equations as backward
SDEs and solve them with a random-feature network, reducing training to a
least-squares regression. They share the forward target but not the mechanism,
and neither dimension scaling nor the inverse problem is studied.

\paragraph{AI methods for the calibration inverse problem}
The dominant approach for fast rough-volatility calibration is the supervised
surrogate of Bayer et al.~\citep{bayer2019deep} and Horvath et al.~\citep{horvath2021deep}: a network trained
offline on precomputed prices or implied volatilities, used as a fast forward map
inside the calibration loop. It is purely data-driven; the network never sees the
governing equation, carries no PDE consistency, and---because it is not derived
from a physics-informed forward solve---cannot be decomposed into a network-error
and a model-error term the way ours can, which is precisely the diagnostic an
inverse-problem user needs to trust or distrust a calibrated fit.

\paragraph{PINN training pathologies and remedies}
PINNs are well-known to fail as forward operators when loss terms sit on
different scales or temporal causality is violated
\citep{krishnapriyan2021characterizing,wang2022when}. The two stabilisers we use
to keep the forward operator convergent and stable---self-adaptive
homoscedastic-uncertainty weights \citep{kendall2018multitask,mcclenny2023self}
and a causal $\tau$-curriculum \citep{wang2024respecting}---directly address
these failure modes.

\paragraph{Foundational methods}
We build on PINNs \citep{raissi2019physics} and deep PDE solvers
\citep{sirignano2018dgm,han2018solving}; on rough volatility
\citep{gatheral2018volatility,eleuch2019characteristic} and its multifactor
Markovian lift \citep{abijaber2019multifactor,abijaber2019lifting}; and on
classical transform pricing \citep{heston1993closed,carr1999option} for our
independent, non-lifted forward-model ground truth.

% ================================================================
\section{Problem formulation}
\label{sec:problem}

\subsection{Rough Heston and the Markovian lift}
Under rough Heston the spot price $S$ and its instantaneous variance $V$ satisfy
\begin{align}
  dS_t &= \sqrt{V_t}\,S_t\,dW^S_t, \\
  V_t  &= V_0 + \int_0^t K(t-s)\,\lambda(\theta - V_s)\,ds
          + \int_0^t K(t-s)\,\nu\sqrt{V_s}\,dW^V_s,
\end{align}
where $W^S$ and $W^V$ are Brownian motions with correlation $\rho$, while
$K(t)=t^{\alpha-1}/\Gamma(\alpha)$ is the fractional kernel with
$\alpha=H+\tfrac12$ and $H<\tfrac12$ \citep{gatheral2018volatility,eleuch2019characteristic}.
The variance is a Volterra process: non-Markovian, non-semimartingale, and
admitting no finite-dimensional forward PDE. The parameter vector
$(H,V_0,\theta,\lambda,\nu,\rho)$ is exactly what the calibration inverse problem
must recover from market data; without a forward PDE there is no forward operator
to embed inside that inverse problem at all. We fix this notation once, because
the paper uses two different subsets of it: the full six-parameter vector is what
the market calibration of Section~\ref{sec:btc} recovers, whereas the parametric
forward operator of Section~\ref{sec:parametric} conditions on the five
parameters $(H,\nu,\rho,\lambda,\theta)$ that enter the PDE coefficients, holding
$V_0$ (a state, not a coefficient) and the maturity fixed at training time.

The Markovian lift replaces $K$ with
$K(t)\approx\sum_{i=1}^{n}c_i e^{-x_i t}$, introducing $n$
Ornstein--Uhlenbeck factors $U^i$,
\begin{equation}
  \begin{split}
  dU^i_t &= \left[-x_i U^i_t + \lambda\left(\theta - V_t\right)\right]dt
           + \nu\sqrt{V^+_t}\,dW^V_t, \qquad U^i_0 = 0,\\
  V_t &= V_0 + \sum_{i=1}^{n} c_i U^i_t,
  \end{split}
  \label{eq:lift}
\end{equation}
where $V^+_t=\max(V_t,0)$
\citep{abijaber2019multifactor,abijaber2019lifting}. Two features of
\eqref{eq:lift} drive everything that follows. First, the factors differ
\emph{only} in their mean-reversion speed $x_i$: they share the same drift
source $\lambda(\theta-V_t)$, the same diffusion coefficient $\nu\sqrt{V^+_t}$
and, critically, the same Brownian increment $dW^V_t$. That common driver is why
the generator in \eqref{eq:pde} below carries a full $(i,j)$ double sum rather
than a diagonal one; the derivation is given in \ref{app:derivation}. Second, the
speeds $(x_i)$ span roughly ten orders of magnitude at small $H$, which is the
source of the stiffness that dictates our numerical schemes
(Section~\ref{sec:eval}). The weights $(c_i,x_i)$ are
placed on a geometric quadrature grid \citep{abijaber2019multifactor}; the
relative $L^2$ kernel error falls monotonically from $25.6\%$ at $n{=}5$ to
$0.39\%$ at $n{=}100$, so raising $n$ genuinely buys forward-model fidelity, at
the cost of raising the forward-problem's dimension by the same $n$.

\subsection{The lifted forward PDE}
The discounted European call price $P(t,x,u)$, where $x=\log S$ and
$u=(u_1,\ldots,u_n)$ collects the factor values, satisfies the
$(n{+}1)$-dimensional backward Kolmogorov PDE
\begin{align}
  0 = {}& P_t + \!\left(r-\tfrac12 V\right)P_x + \tfrac12 V\,P_{xx}
         + \sum_{i=1}^{n}\!\left[-x_i u_i + \lambda(\theta-V)\right]P_{u_i}
         \nonumber\\
        & + \tfrac12\nu^2 V\sum_{i,j=1}^{n}P_{u_i u_j}
         + \rho\nu V\sum_{i=1}^{n} P_{x u_i} - rP ,
  \label{eq:pde}
\end{align}
with terminal condition $P(T,x,u)=\max(e^x - K, 0)$. Because all factors share
the single Brownian motion $W^V$, the full double sum over $(i,j)$ and the
$(n{+}1)$-dimensional cross term both appear; the PDE dimension is exactly
$n{+}1$. At $n{=}4$ this is a five-dimensional forward problem; at $n{=}32$ it is
thirty-three. Solving \eqref{eq:pde} for a single candidate parameter vector is
one forward-model evaluation inside the outer calibration loop; a naive
optimiser needs hundreds to thousands of these per fit (Section~\ref{sec:calib}).

% ================================================================
\section{The PINN forward operator: architecture and training}
\label{sec:arch}

\subsection{Baseline network}
We represent the forward map as a fully connected multilayer perceptron (MLP).
Every hidden unit uses the $\tanh$ activation, which provides the smooth,
high-order derivatives needed for automatic differentiation of the PDE residual;
$\tanh$ is standard for PINN residual training and avoids the dead-neuron
pathology of ReLU in this high-derivative setting \citep{raissi2019physics}. All
weight matrices are initialised with Xavier normal initialisation
\citep{glorot2010understanding} and all biases with zeros.

The network maps $(\tau, x, u_1,\ldots,u_n) \in \mathbb{R}^{2+n}$ to a scalar
correction. For the single-point forward operator (parameters fixed) we use
width $W{=}64$ and depth $D{=}4$ hidden layers; for the parametric forward
operator used inside the inverse problem (Section~\ref{sec:parametric}) we
increase to $W{=}96$, $D{=}5$. Table~\ref{tab:arch} summarises all architectural
and training hyperparameters.

\begin{table}[!htbp]
\centering
\caption{Architecture and training hyperparameters of the PINN forward operator.
Values in parentheses are for the parametric extension (Section~\ref{sec:parametric}).}
\label{tab:arch}
\begin{tabularx}{\linewidth}{@{}l >{\raggedright\arraybackslash}X@{}}
\toprule
\textbf{Component} & \textbf{Value} \\
\midrule
\multicolumn{2}{l}{\textit{Architecture}} \\
Network type & Fully connected MLP \\
Hidden layers & 4 (5 for parametric) \\
Neurons per layer & 64 (96 for parametric) \\
Activation & $\tanh$ (all hidden layers) \\
Weight initialisation & Xavier normal \\
Bias initialisation & Zero \\
Input dimension & $2 + n$ for single-point; $2 + n + 5$ for parametric \\
Output dimension & 1 (price correction $N_\theta$) \\
Floating-point precision & \texttt{float64} (PyTorch) \\
\midrule
\multicolumn{2}{l}{\textit{Optimisation}} \\
Optimiser & Adam \citep{kingma2015adam} \\
Initial learning rate & $10^{-3}$ \\
LR schedule & Cosine annealing to $0$ over all iterations \\
Gradient clipping & $\ell_2$ norm $\leq 1.0$ \\
Training iterations & 10,000--20,000 (single-point); 20,000--30,000 (parametric) \\
\midrule
\multicolumn{2}{l}{\textit{Collocation sampling (per iteration)}} \\
Interior collocation & 2,000 uniform + 500 near-strike (total 2,500) \\
Boundary points & 400 (split OTM / ITM) \\
Log-moneyness domain & $[\log K - 1.2,\;\log K + 1.2]$ \\
Factor sampling & MC-empirical mean / std (Section~\ref{sec:mcsamp}) \\
Tail coverage fraction & 20\% of $u$-points drawn at $2\times$ std \\
\midrule
\multicolumn{2}{l}{\textit{Loss weighting}} \\
Adaptive scheme & Learnable log-variance (homoscedastic uncertainty) \\
Terms balanced & PDE residual MSE, OTM boundary MSE, ITM boundary MSE \\
\midrule
\multicolumn{2}{l}{\textit{$\tau$-curriculum}} \\
Initial horizon & $\tau_0 = 0.15\,T$ \\
Ramp fraction & Linear growth to $T$ over first $50\%$ of iterations \\
\midrule
\multicolumn{2}{l}{\textit{Compute environment (every run in this paper)}} \\
\IfFileExists{tables/environment.tex}{\revtabletrue}{\revtablefalse}\ifrevtable
\revrows{environment}
\else
Platform & CPU (single machine); code is GPU-compatible (\texttt{torch.device}) \\
\fi
\bottomrule
\end{tabularx}
\end{table}

\subsection{Hard-constrained initial-condition ansatz}
\label{sec:ansatz}
We train in time-to-maturity $\tau = T-t$ so the terminal condition becomes an
initial condition at $\tau{=}0$. Rather than penalise it softly, we enforce it
exactly through the ansatz
\begin{equation}
  P(\tau,x,u) = P_{\mathrm{base}}(\tau,x) + \tau\,N_\theta(\tau,x,u),
  \label{eq:ansatz}
\end{equation}
where $N_\theta$ is the network output. The $\tau$ prefactor forces the learned
correction to vanish at $\tau{=}0$ automatically, so the network never fights the
initial condition. Without this, a soft penalty for the terminal payoff
oscillates against the PDE residual and the forward operator fails to converge.

\subsection{Integrated-variance baseline anchor}
\label{sec:baseline}
The baseline $P_{\mathrm{base}}$ is a Black--Scholes call price. A critical
design choice is what volatility it is anchored at. Anchoring at the
\emph{spot} variance $V_0$ leaves a systematic level bias: the factors
mean-revert from $V_0$ toward $\theta$ over the life of the option, so the
model's actual expected variance is higher than $V_0$. We anchor instead at the
model's expected \emph{integrated} variance
\begin{equation}
  w(\tau) = \int_0^\tau \mathbb{E}[V_s]\,ds,
  \label{eq:intvar}
\end{equation}
computed by a short pre-flight Monte Carlo (the same simulation used for factor
collocation). The baseline then uses $\sigma(\tau) = \sqrt{w(\tau)/\tau}$, which
correctly captures the mean-reversion effect; the network learns only the
residual stochastic-volatility correction. This single change removed a uniform
${\sim}120$ vol\bp{} level bias we observed with the $V_0$ anchor---an error that
would otherwise propagate directly into every inverse-problem fit using this
forward operator.

\subsection{MC-informed factor collocation}
\label{sec:mcsamp}
The factor values $u$ at a given maturity $\tau$ are not centred at zero: they
mean-revert to a distribution set by the model dynamics. Sampling from a fixed
symmetric box centred at zero causes the network to expend capacity on regions
the dynamics never visit, and---critically---to under-cover the region governing
the put-wing skew. Before training we run a short lifted Monte Carlo simulation
(4,000 paths, 300 time steps) to estimate the mean and standard deviation of each
factor at the target maturity; training collocation is then drawn from
$\mathcal{N}(\hat{\mu}_i, \hat{\sigma}_i)$. An additional 20\% of collocation
points are drawn at twice the estimated standard deviation to extend tail
coverage. This change cut smile RMSE by roughly $45\%$ (from ${\sim}210$ to
${\sim}114$ vol\bp) relative to symmetric-box sampling.

\subsection{Self-adaptive loss weighting}
\label{sec:adaptive}
Three loss terms must be balanced at each training step: the PDE-residual MSE
$\mathcal{L}_{\mathrm{PDE}}$, and the MSEs at the deep-OTM ($P\to0$) and
deep-ITM ($P\to S - Ke^{-r\tau}$) boundaries. Fixed-weight balancing is brittle:
the ITM boundary target is $O(1)$ while the OTM target is $O(0)$, so a single
lumped weight over-emphasises the ITM term. We adopt the learnable
homoscedastic-uncertainty scheme of Kendall et al.~\citep{kendall2018multitask} in the
self-adaptive spirit of McClenny and Braga-Neto~\citep{mcclenny2023self}: each term $\mathcal{L}_k$ is
assigned a learnable log-variance $s_k$ and the objective is
\begin{equation}
  \mathcal{L} = \sum_{k\in\{\mathrm{PDE},\,\mathrm{OTM},\,\mathrm{ITM}\}}
               \bigl[e^{-s_k}\,\mathcal{L}_k + s_k\bigr].
  \label{eq:adaptive}
\end{equation}
The $s_k$ are optimised jointly with the network weights; large or noisy terms
are down-weighted automatically. We never hand-tune a boundary weight.

\subsection{Causal \texorpdfstring{$\tau$}{tau}-curriculum}
\label{sec:curriculum}
Forward operators that train simultaneously across all time-to-maturities can be
pulled toward spurious solutions at long $\tau$ before the short-$\tau$ region
has converged, because the residual at long $\tau$ depends causally on the
short-$\tau$ solution \citep{wang2024respecting}. We address this with a causal
curriculum: training begins at $\tau_{\mathrm{max}} = 0.15\,T$ and linearly ramps
to the full $T$ over the first half of training. The short-$\tau$ solution is
therefore established before the long-$\tau$ residual can perturb it.

\subsection{Training algorithm}
Algorithm~\ref{alg:train} summarises the full training procedure for the forward
operator.

\begin{algorithm}[t]
\caption{Training the PINN forward operator for the lifted rough-Heston PDE}
\label{alg:train}
\begin{algorithmic}[1]
\Require Model parameters $(V_0,\lambda,\theta,\nu,\rho)$; lift $(c_i,x_i)_{i=1}^n$;
         maturity $T$; hyperparameters in Table~\ref{tab:arch}
\State Initialise MLP $N_\theta$ (Xavier normal); learnable $s = [s_1,s_2,s_3] \leftarrow 0$
\State \textbf{Pre-flight MC:} simulate lifted SDE, record $(\hat\mu_i, \hat\sigma_i)$
       for each factor and term structure $w(\tau)$ \Comment{Sections~\ref{sec:mcsamp},~\ref{sec:baseline}}
\For{iteration $t = 1, \ldots, N_{\mathrm{iter}}$}
  \State $\tau_{\mathrm{max}} \leftarrow \min\!\bigl(T,\; 0.15T + \tfrac{0.85T \cdot t}{0.5\,N_{\mathrm{iter}}}\bigr)$
         \Comment{$\tau$-curriculum, Sec.~\ref{sec:curriculum}}
  \State Sample interior: $\{(\tau_j, x_j, u_j)\}$ from
         $\mathrm{Unif}([0,\tau_{\mathrm{max}}])\times\mathrm{Unif}(x_{\mathrm{lo}},x_{\mathrm{hi}})
         \times\mathcal{N}(\hat\mu,\hat\sigma)$
  \State Sample boundary: $\{(\tau^{\mathrm{bc}}_j, x_{\mathrm{lo/hi}}, u^{\mathrm{bc}}_j)\}$
  \State $\hat P \leftarrow P_{\mathrm{base}}(\tau,x;\,w) + \tau\,N_\theta(\tau,x,u)$
         \Comment{Ansatz, Eq.~\eqref{eq:ansatz}}
  \State $\mathcal{L}_{\mathrm{PDE}} \leftarrow$ mean-squared PDE residual of $\hat P$
         (Eq.~\eqref{eq:pde})
  \State $\mathcal{L}_{\mathrm{OTM}} \leftarrow \|\hat P(\tau^{\mathrm{bc}},x_{\mathrm{lo}},u^{\mathrm{bc}})\|^2$,\quad
         $\mathcal{L}_{\mathrm{ITM}} \leftarrow \|\hat P(\tau^{\mathrm{bc}},x_{\mathrm{hi}},u^{\mathrm{bc}}) - (e^{x_{\mathrm{hi}}} - Ke^{-r\tau^{\mathrm{bc}}})\|^2$
  \State $\mathcal{L} \leftarrow \sum_k [e^{-s_k}\mathcal{L}_k + s_k]$
         \Comment{Adaptive weighting, Eq.~\eqref{eq:adaptive}}
  \State Backpropagate $\mathcal{L}$; clip gradients ($\ell_2 \leq 1.0$); Adam step
         on $(N_\theta, s)$; cosine LR anneal
\EndFor
\end{algorithmic}
\end{algorithm}

% ================================================================
\section{Evaluation methodology: trusting the forward operator}
\label{sec:eval}

\subsection{Two independent ground truths}
An AI forward operator used inside an inverse problem is only as trustworthy as
its own error against the true forward map. We hold the network to two
independent references.

\textbf{Fourier forward model.} A semi-analytic pricer that solves the
fractional Riccati equation for the rough Heston characteristic function
\citep{eleuch2019characteristic} and inverts by Fourier integration
\citep{carr1999option}. It never uses the lift, so it is the ground truth of the
\emph{true, non-lifted model} ($n\to\infty$); in the flat-volatility limit it
reproduces Black--Scholes to $3.7\times10^{-7}$ in price.

\textbf{Lifted Monte Carlo.} An exponential-integrator Euler scheme that
integrates the $n$-factor lifted SDE with the $-x_i U^i$ drift handled exactly:
\begin{equation}
  U^i_{t+\Delta t} = e^{-x_i\Delta t}\,U^i_t
    + \int_t^{t+\Delta t}e^{-x_i(t+\Delta t - s)}\sqrt{V_s}\,dW^V_s.
  \label{eq:expint}
\end{equation}
The exponential-integrator is essential: the mean reversions $(x_i)$ span roughly
ten orders of magnitude for small $H$, and a naive explicit Euler scheme applied
to the factor equations diverges. The lifted MC matches Black--Scholes within
Monte Carlo error in the flat-variance degeneracy ($\nu{=}\lambda{=}0$;
100,000 paths, 200 steps) and agrees with the Fourier forward model for large
$n$. At finite $n$, the lifted MC is the \emph{exact solution of the forward PDE
the network is asked to solve}; the Fourier forward model is the exact solution
of the true model.

\paragraph{Precision of the reference} Since every solve-error in this paper is
measured against this simulator, its own error budget has to be smaller than the
quantity being measured, and we state it explicitly rather than leave it
implicit. All validation runs use $10^{6}$ antithetic paths and $1\,600$ time
steps with a terminal-forward control variate. The statistical standard error is
$0.69$ px\bp{} at $n{=}4$ rising to $0.81$ px\bp{} at $n{=}32$, against
solve-errors of $15$--$22$ px\bp{}. Two points deserve emphasis. First, the
standard error must be computed on the antithetic \emph{pair means}, not on the
individual paths; the naive independent-sample formula is not a valid estimator
here and overstates the RMS error by $14$--$31$ per cent. Second, the
time-discretisation bias is a systematic error that no amount of sampling
removes, and it is the more dangerous of the two. A common-random-numbers
refinement study (Table~\ref{tab:mcq}) measures it at $0.21$, $0.17$, $0.82$ and
$1.27$ px\bp{} for $n{=}4,8,16,32$ at $1\,600$ steps. We record why the step
count is $1\,600$ and not the $400$ of the submitted version, because the choice
was forced by the data. At the calibration of Section~\ref{sec:btc} the factor
system is considerably stiffer than at the parameters we first validated on
($\nu=1.37$ and $\lambda=4.94$ against $0.30$ and $2.57$), and at $400$ steps the
bias grows to $3.2$ and $3.8$ px\bp{} at $n{=}16$ and $32$. That is comparable
to the lift-error being measured, and it was visible: against the $400$-step
reference the lift-error at those $n$ stopped falling and changed sign. At
$1\,600$ steps the bias is below a tenth of every solve-error in this paper.

\IfFileExists{tables/mcquality.tex}{\revtabletrue}{\revtablefalse}\ifrevtable
\begin{table}[!htbp]
\centering
\caption{Precision of the lifted Monte Carlo reference at the calibration of
Section~\ref{sec:btc}, $T{=}0.15$, in price basis points. Standard errors are
RMS over strikes; the naive column is the independent-sample formula, which is
not valid under antithetic sampling and is shown only to quantify the
overstatement. Time-discretisation bias is from a common-random-numbers
refinement against $3\,200$ steps.}
\label{tab:mcq}
\small
\begin{tabular}{@{}r r r r r r@{}}
\toprule
$n$ & SE, $2{\times}10^5$ & naive SE & SE, $10^6$ & bias, 400 & bias, 1\,600 \\
\midrule
\revrows{mcquality}
\bottomrule
\end{tabular}
\end{table}
\fi

\subsection{Error decomposition}
\label{sec:decomp}
Rather than reporting a single ``network-versus-true-model'' figure, we decompose
the error per strike and in price space (where the signed contributions add
exactly),
\begin{equation}
  \underbrace{(\mathrm{PINN}-\mathrm{Fourier})}_{\text{total forward error}}
  =\underbrace{(\mathrm{PINN}-\mathrm{MC})}_{\text{network solve-error}}
  +\underbrace{(\mathrm{MC}-\mathrm{Fourier})}_{\text{lift (model) error}}.
  \label{eq:split}
\end{equation}
The \emph{solve-error} measures how accurately the network solves the forward
PDE it was given; it is the quantity training controls, and it is the AI method's
own contribution to inverse-problem error. The \emph{lift-error} is a
forward-model approximation that only increasing $n$ can reduce; it is not the
network's responsibility, and no amount of network re-training removes it.
Conflating the two gives a systematically pessimistic view of the AI method and
leads to the wrong engineering response inside a calibration pipeline---e.g.\ training longer when the correct response is to raise $n$.

\paragraph{Units, stated once} Two units appear in this paper and a referee
rightly asked us to connect them explicitly. Price errors are in basis points of
spot (px\bp; $1\,\mathrm{px\,bp}=10^{-4}\times S$), and we use these wherever
the decomposition \eqref{eq:split} is invoked, because it is additive in price
and \emph{not} in implied volatility. Implied-volatility errors are in vol\bp{}
($1\,\mathrm{vol\,bp}=10^{-4}$ in volatility), which is the unit a practitioner
reads. The two are related by the Black--Scholes vega: at the money, $T{=}0.15$
and $S{=}1$, vega is $0.154$ almost independently of the volatility level, so
\[
  100\,\mathrm{vol\bp} \;\approx\; 15.4\,\mathrm{px\bp}
  \qquad\text{(at the money, }T{=}0.15\text{)}.
\]
The conversion is pointwise, not global: vega falls away in the wings, so an
RMSE aggregated over the strike range $0.85$--$1.15$ converts at roughly
$12$--$14$ px\bp{} per $100$ vol\bp{} rather than $15.4$. This is why, for
instance, a lift-error of $16.4$ px\bp{} in Table~\ref{tab:nconv} corresponds to
$133$ vol\bp{} in Section~\ref{sec:decomp_results} rather than to $106$. To keep
the two readable side by side, \textbf{Table~\ref{tab:nconv} now reports both
units for every entry}. Finally, note that the two components of
\eqref{eq:split} carry opposite signs, so their RMSE values do not add even
though the signed errors do.

\subsection{Homogeneity re-scaling}
We train the network at fixed strike $K{=}1$ and read off the smile $C(1,K)$
through the call's degree-one homogeneity:
\begin{equation}
  C(1,K) = K\,C\!\left(1/K,\;1\right),
  \label{eq:homog}
\end{equation}
which holds exactly because the variance dynamics are independent of the spot
level. Dropping this re-scaling and comparing mismatched options inflates the
apparent error to ${\sim}2700$ vol\bp; we verified this is a comparison artefact,
not a modelling failure of the forward operator.

% ================================================================
\section{The calibration inverse problem: data and methods}
\label{sec:calib}

\subsection{Data}
\label{sec:data}
We use a single BTC option-chain snapshot from Deribit, a cryptocurrency
derivatives exchange, retrieved through its public API (no API key required) on
\textbf{3 July 2026 at 05:57 UTC}, at a spot of \$62{,}191. The raw snapshot
holds 912 quoted instruments across 13 expiries.

Because the referees asked for the data handling to be stated precisely, we give
the filter chain in full. From the raw chain we (i) take the mid price
$\tfrac12(\text{bid}+\text{ask})$ where a two-sided quote exists and otherwise
fall back to Deribit's \emph{mark} price; (ii) keep only the out-of-the-money
leg at each strike (calls for $K\ge F$, puts for $K<F$), which is the more
liquid and tighter-quoted side; (iii) drop $|k|>0.30$ with $k=\log(K/F)$;
(iv) drop quotes below one currency unit, where the implied vol is dominated by
tick noise; and (v) recompute the implied volatility from the mid by Brent
inversion, dropping any point falling outside $[0.01,2.0]$. This leaves 286
points, from which we retain at most eight per maturity, spread evenly in
log-moneyness so that each smile constrains shape and not merely level. The
fitted surface therefore comprises \textbf{104 points across 13 maturities}
($T$ from $0.002$ to $0.98$ years) spanning a moneyness range
$K/F\in[0.77,1.35]$; of these, 97 return a finite model implied volatility at
the calibrated parameters and enter the reported RMSE.

Two disclosures about that surface. First, \textbf{11 of the 104 fitted points
carry no two-sided quote} and therefore enter through Deribit's mark price,
which is an exchange model value rather than a tradable one. Second,
\textbf{12 of the 104 had no traded volume} in the snapshot window. We do not
filter on volume or open interest: Deribit's book-summary endpoint does not
return open interest, and imposing a volume filter would remove most of the long
maturities entirely. Section~\ref{sec:btc} reports the fit restricted to the
liquid moneyness core as a robustness check, and
Section~\ref{sec:arbitrage} reports static-arbitrage diagnostics on the surface
that survives the filters above.

\paragraph{Two further datasets, used only for robustness}
The regime and roughness checks of Section~\ref{sec:roughness_robustness} use
two datasets distinct from the headline snapshot, and we describe both here so
that their caveats sit with the data rather than with the results. (i) A
five-minute bar history of the Deribit BTC perpetual, August 2018 to August
2026 ($841$k bars), collected by the authors through the same public API. It
enters only the realised-volatility estimate of $H$ and never touches an
option. (ii) A daily implied-volatility panel of BTC options, December 2016 to
August 2026, built by the authors from the Deribit public trade tape ($41.5$M
option trades). Every entry in that panel is a completed transaction, not a
quote: there is no bid--ask, the two-sided-quote filter above cannot be
applied, and a single stale print can set the implied volatility of an
illiquid strike. We therefore restrict the panel to $|k|\le0.35$,
$T\in[0.02,1.0]$ years and implied volatilities in $[0.01,2.0]$, retain at
most eight points per maturity spread in log-moneyness exactly as above, and
keep only Wednesdays with at least four maturities and 25 points. Trade-based
and quote-based surfaces are each internally consistent but are not
interchangeable: a comparison across weeks of the panel is sound, but the
panel is not a like-for-like extension of the headline snapshot, and we do
not present it as one. Both datasets accompany the code in derived form (the
weekly surfaces and the daily realised variance), not as raw tapes.

We set $r{=}0$, the standard convention for crypto, where there is no clean
carry, and take the forward to be $F=S\,e^{rT}=S$ with $S$ the underlying index
price Deribit reports alongside each instrument. Implied volatilities are
computed from mid prices with a robust Brent Black--Scholes inverter. The
inverse-problem objective throughout is the vega-weighted implied-volatility
RMSE between market quotes and the forward operator's output,
\begin{equation}
  \mathcal{L}(\Theta)=\sqrt{\frac{\sum_j w_j\left[\sigma^{\mathrm{mod}}_j(\Theta)
        -\sigma^{\mathrm{mkt}}_j\right]^2}{\sum_j w_j}},
  \qquad
  w_j=\max\!\left(S\,\phi(d_{1,j})\sqrt{T_j},\;10^{-6}\right),
  \label{eq:objective}
\end{equation}
with $\phi$ the standard normal density and $d_{1,j}$ evaluated at the market
implied volatility, so the weights are fixed across the optimisation. Vega
weighting lets well-determined near-the-money points dominate, which is standard
practice. We flag one asymmetry for the reader: the \emph{objective} is
vega-weighted as in \eqref{eq:objective}, whereas the headline fit quality we
report in Section~\ref{sec:btc} is the \emph{unweighted} RMSE over the same
points, which is the more conservative of the two.

\subsection{Static-arbitrage diagnostics}
\label{sec:arbitrage}
Attributing a coarse fit to the forward model's expressiveness, as we do in
Section~\ref{sec:btc}, is only legitimate if the target surface is itself
arbitrage-consistent: quotes carrying butterfly or calendar violations cannot be
fitted by \emph{any} arbitrage-free model, and the residual would then be the
data's doing rather than rough Heston's. Crypto surfaces are wide and thinly
quoted in the wings, so this is a real possibility and we test it rather than
assume it.

On undiscounted calls with $r{=}0$, putting the out-of-the-money puts on a
common footing through parity, we check three conditions on the filtered
surface: convexity in strike at each maturity (butterfly), the slope bound
$-1\le \partial C/\partial K\le 0$ (vertical spread), and monotonicity of total
implied variance $w(k,T)=\sigma^2 T$ in $T$ at fixed log-moneyness (calendar).
The result is clean: \textbf{no violations} in 78 butterfly triples, 91 vertical
pairs or 300 calendar comparisons. The margins are however slim --- the smallest
butterfly margin is $0.28$ px\bp{} and the smallest calendar increment is
$8.2\times10^{-5}$ in total-variance units --- so the surface is arbitrage-free
but only barely, which is itself a fair description of a wide crypto smile. The
calibrated model's own surface passes the same tests with larger margins
($0.98$ px\bp{} butterfly), as an arbitrage-free model must.

Separately, put--call parity holds on the raw chain: of 383 strikes quoted
two-sided on both legs, \textbf{none} violates $C-P=F-K$ by more than the quoted
bid--ask spread, though the residual RMS of $46$ px\bp{} (maximum $163$) is a
reminder of how wide those spreads are.

The conclusion we draw is narrow but useful: the coarseness of the fit reported
in Section~\ref{sec:btc} \emph{cannot} be blamed on static-arbitrage violations
in the target, because there are none. It is a statement about the forward model.
\IfFileExists{tables/liquid.tex}{\revtabletrue}{\revtablefalse}\ifrevtable

\paragraph{The liquid core}
The sharper form of the referee's question is whether the coarse fit is driven by the wings.
Table~\ref{tab:liquid} refits on the liquid core, $|k|\le0.10$ and $|k|\le0.20$, from both
ends of the $H$ ridge of Section~\ref{sec:roughness_robustness}, and also evaluates the two
unrefitted full-surface vectors on the same core points, which separates the effect of
dropping the wings from the effect of re-optimising.
Two things follow from Table~\ref{tab:liquid}, and the comparison that matters is
on a common set of points. First, the wings are not what makes the fit coarse.
On the fifty-three points with $|k|\le0.10$, the full-surface fit scores
$160.4$ vol\bp{} weighted and $510.8$ unweighted; refitting on those points alone
improves this only to $146.4$ and $450.0$, a gain of nine to twelve per cent. A
model that remains above $450$ vol\bp{} on the most liquid quotes it is shown is
not being defeated by illiquid ones; the residual belongs to the lifted
rough-Heston form. Second, and more interesting, the recovered $H$ depends
strongly on the window: $0.020$ on $|k|\le0.10$, $0.057$ on $|k|\le0.20$, and
$0.172$ on the full surface. On both restricted windows the two warm starts
converge to the same parameters---to $0.2$ vol\bp{} on the core and to the
reported precision on $|k|\le0.20$---so there $H$ \emph{is} identified. The
non-identification of Section~\ref{sec:roughness_robustness} is therefore
specifically a property of the far wings, which is where the smoother-kernel,
higher-vol-of-vol end of the ridge buys its fit.
\begin{table}[!htbp]
\centering
\caption{Liquid-core refit. Fit quality in vol\bp, weighted by the vega weights of the objective and unweighted. Unrefitted rows evaluate the full-surface vectors on the restricted points.}
\label{tab:liquid}
\footnotesize
\begin{tabular}{@{}l r l r r r r r@{}}
\toprule
Window & pts & Fit & weighted & unweighted & $H$ & $\rho$ & $\nu$ \\
\midrule
\revrows{liquid}
\bottomrule
\end{tabular}
\end{table}
\fi

\subsection{Three routes to the inverse problem}
We compare three ways of solving the same calibration inverse problem.
\begin{enumerate}
  \item \textbf{True-forward-model loop.} Differential evolution (60
        generations, all 5 parameters) with the Fourier forward model inside the
        objective, followed by Nelder--Mead polish. This is our from-scratch,
        exact-forward-operator reference.
  \item \textbf{AI-only.} Gradient descent on the model parameters using the
        parametric PINN as a differentiable forward operator, with zero calls to
        the true forward model.
  \item \textbf{Hybrid.} The parametric PINN's single-shot solution provides the
        initial guess; a short Nelder--Mead polish against the true forward
        model then sharpens it.
\end{enumerate}

\paragraph{Optimiser settings, stated once}\label{par:optimiser}
Two distinct budgets appear in this paper and conflating them would misstate the
speed-up, so we separate them explicitly. The \emph{full-budget market
calibration} of Section~\ref{sec:btc} uses SciPy's \texttt{differential\_evolution}
with \texttt{maxiter}${=}60$, \texttt{popsize}${=}12$, \texttt{tol}${=}10^{-4}$,
mutation $(0.5,1.0)$, recombination $0.7$ and seed $1$, the built-in polish
disabled, evaluated in parallel with deferred population updating, followed by a
separate Nelder--Mead refinement (\texttt{maxiter}${=}200$,
\texttt{xatol}${=}10^{-4}$, \texttt{fatol}${=}10^{-6}$) confined to the box. The
box is $H\in[0.02,0.45]$, $V_0,\theta\in[0.005,0.60]$, $\lambda\in[0.05,12]$,
$\nu\in[0.05,2]$, $\rho\in[-0.95,-0.05]$; every parameter of the reported fit is
interior to it. The \emph{three-route comparison} of this section runs on the
single $T{=}0.15$ slice with a deliberately lighter global budget---\texttt{maxiter}${=}30$,
\texttt{popsize}${=}10$, seed $1$---and a Nelder--Mead polish capped at $120$
iterations. All forward-model evaluation counts we report, including the
$1\,736$-evaluation from-scratch baseline against which the $8.8\times$ saving is
measured, come from this lighter configuration; they are not comparable to the
full-budget market fit, and we do not claim they are.

We also record how the box was arrived at, because it was tight three times
over and each time it moved the answer. In the submitted fit the Nelder--Mead
refinement ran \emph{unbounded} and $\theta$ escaped the interval the global
search had explored. Under the corrected objective, a box capping $\lambda$ at $3$
left $\lambda$ on its bound; the box above is wide enough to contain every
parameter. In the profile likelihood of Section~\ref{sec:roughness_robustness}
the same cap held $\lambda$ on its bound at ten of twelve grid points and narrowed
the reported $H$ interval by almost half until it was lifted. A search box that
constrains rather than contains the optimum is a quiet source of false
precision, and we now check every reported fit for bound contacts.

% ================================================================
\section{Results}
\label{sec:results}

\subsection{The forward operator scales where the grid cannot}
\label{sec:nconv}
Our headline forward-operator experiment trains one architecture (width 64,
depth 4) at $n{=}4,8,16,32$ on the BTC calibration of Section~\ref{sec:btc} at
$T{=}0.15$, over eight independent seeds per $n$ (Table~\ref{tab:nconv},
Figure~\ref{fig:nconv}). The solve-error against the lifted MC is $15.2\pm4.1$,
$21.2\pm5.2$, $24.6\pm5.0$ and $18.3\pm6.1$ px\bp{} at $n{=}4,8,16,32$. It rises
to $n{=}16$ and falls back at $n{=}32$, and across the eight-fold rise in
dimension it grows by a factor of $1.21$. We test the question the scaling claim
turns on directly, on all thirty-two seed-level points. A least-squares fit of
log solve-error on $\log n$ gives a slope of $0.09\pm0.08$, significantly below
one, so the growth is sub-linear; and a one-sided permutation test of
monotonically increasing error, over $20\,000$ relabellings, gives $p=0.14$, so
we cannot reject that the solve-error is flat in $n$. The submitted version
reported a doubling on three seeds whose error bars overlapped at adjacent $n$;
with eight seeds the defensible statement is the stronger one, that we detect no
growth at all. The absolute level is higher than in the submitted version
because every experiment now runs at the corrected calibration, whose fast
mean-reversion ($\lambda=4.94$) and large vol-of-vol ($\nu=1.37$) make the
lifted factor system markedly stiffer. Training cost does not grow with $n$
either: measured on an otherwise idle machine, one training iteration costs
$84$--$120$\,ms at every $n$, with no trend beyond the scatter between repeats,
because the factor-Hessian computation is vectorised over the full double sum.

The lift-error falls from $8.8$ px\bp{} at $n{=}4$ to $3.8$ at $n{=}8$ and reaches
the resolution floor of the reference by $n{=}16$ (shaded in
Figure~\ref{fig:nconv}): the Monte Carlo reference's standard error and
time-discretisation bias together amount to $1.8$ px\bp{} at $n{=}16$ and $2.7$
at $n{=}32$, and the lift-errors measured there, $1.0$ and $3.4$, are not
resolvable against it. Two checks show the floor is the reference and not the
Fourier price: the Fourier price is converged to $0.1$ vol\bp{} across a
sixteen-fold refinement of both its time grid and its frequency truncation, and
the signed lift-error rises monotonically in $n$ at every strike, crossing zero
between $n{=}8$ and $16$ at and above the money and between $16$ and $32$ in the
put wing. That is the signature of a kernel-approximation error that vanishes
with $n$ meeting a positive discretisation bias that grows with it, which is
what the floor itself shows. We report the lift-error beyond $n{=}16$ as
below the floor and draw no conclusion from its value there.

Compare this against a finite-difference forward solver. A naive full-tensor
grid on the same $(n{+}1)$-dimensional PDE at a modest 50 nodes per dimension
needs $50^{n+1}$ points: $3.1\times10^{8}$ at $n{=}4$, $2.0\times10^{15}$ at
$n{=}8$, and $1.2\times10^{56}$ at $n{=}32$. Even $n{=}8$---roughly $16$
petabytes in double precision---is out of reach as a forward operator inside any
iterative inverse-problem solver, while the network trains every $n$ on a
laptop. We stress that $50^{n+1}$ is an illustrative memory bound for the naive
full-tensor discretisation and not a benchmark against a tuned solver.
Structured schemes reach considerably further: sparse grids
\citep{bungartz2004sparse,reisinger2007efficient} break the full-tensor node
count, ADI splitting handles the correlation cross-terms efficiently in the
classical Heston setting \citep{hout2010adi}, and low-rank tensor-train
representations \citep{oseledets2011tensor} compress high-dimensional operators
where the solution admits low-rank structure. Each, however, still degrades with
$n$---sparse-grid node counts and error constants grow super-linearly in the
dimension, and the tensor-train ranks that make the method cheap are eroded by
precisely the coupled $\rho$--$\nu$ cross-terms this PDE carries---and none is
standard or validated for the lifted rough-volatility PDE. We therefore make the
conservative claim, and only it: a competitive grid-based forward solver at
$n{=}16$--$32$ is \emph{impractical} as the engine of a calibration loop, not
provably impossible. Establishing the stronger claim would require building and
tuning such a solver, which is a research programme in its own right.

\begin{table}[!htbp]
\centering
\caption{Dimension scaling of the forward operator: single-point PINN at
$T{=}0.15$ on the BTC calibration of Section~\ref{sec:btc}. Solve-error is the
eight-seed mean\,$\pm$\,SD in price basis points of spot, against a lifted Monte
Carlo of $10^6$ antithetic paths at $1\,600$ steps; lift-errors at $n{=}16$ and
$32$ are below that reference's resolution floor (see text). Cost is milliseconds
per training iteration, measured with no other job running.}
\label{tab:nconv}
\small
\begin{tabular}{@{}r r r r r@{}}
\toprule
$n$ & solve-error (px\bp) & lift-error (px\bp) & ms / iteration & FD grid $50^{(n+1)}$\\
\midrule
%% rows come from make_tables.py (results/n_convergence.json, results/n_cost.json);
%% the hand-typed submitted-version rows below are a fallback only
\IfFileExists{tables/nconv.tex}{\revtabletrue}{\revtablefalse}\ifrevtable
\revrows{nconv}
\else
4  & $5.4\pm1.0$  & $16.4$ & $397$ & $3.1\times10^{8}$  \\
8  & $6.1\pm1.4$  & $7.9$  & $386$ & $2.0\times10^{15}$ \\
16 & $8.8\pm1.6$  & $3.5$  & $591$ & $7.6\times10^{28}$ \\
32 & $10.7\pm2.0$ & $1.7$  & $457$ & $1.2\times10^{56}$ \\
\fi
\bottomrule
\end{tabular}
\end{table}

\begin{figure}[!htbp]\centering
\includegraphics[width=0.82\linewidth]{fig1_n_convergence.png}
\caption{Dimension scaling of the forward operator. PINN solve-error (blue,
eight-seed mean\,$\pm$\,SD) shows no statistically detectable growth over an
eight-fold rise in $n$ (log-log slope $0.09\pm0.08$; permutation $p=0.14$), and
training cost per iteration (green) is flat. The lift-error (red) falls into the
Monte Carlo reference's resolution floor (grey: standard error combined with
time-discretisation bias) by $n{=}16$, beyond which it is not resolvable. The
inset shows the naive finite-difference grid size $50^{n+1}$, infeasible by
$n{=}8$.}
\label{fig:nconv}
\end{figure}

\IfFileExists{tables/hparams_grid.tex}{\revtabletrue}{\revtablefalse}\ifrevtable
\subsection{Sensitivity to architecture and collocation, and where the variation lives}
\label{sec:hparams}
Two questions follow from Table~\ref{tab:nconv}. Is the width-64 / depth-4 choice a
knife-edge? And the solve-error, though it shows no significant trend in $n$, does
rise and fall across $n$ by up to sixty per cent: does that variation come from the
lifted PDE or from the optimiser? Table~\ref{tab:hpgrid} answers the first with a $3\times3$ architecture grid
at $n{=}4$. For the second we use the a-priori structure of PINN error estimates
\citep{mishra2023estimates,deryck2024error}: the solve-error is bounded by a stability
constant of the PDE times the residual, so we score every trained network on a fixed,
seeded, never-trained-on set of $20\,000$ collocation points and follow that held-out
residual across $n$ (Table~\ref{tab:hpnsweep}). If the residual is flat while the
solve-error moves, the variation sits in the stability constant $C(n)$; if the residual
moves with it, the cause is optimisation or quadrature, and Table~\ref{tab:hpcolloc} tests the
quadrature half directly by raising the collocation count at $n{=}32$, where the paper holds
$n_{\mathrm{col}}$ at $2000$ while the dimension is $33$.
Three findings, the first two against us. The architecture used throughout this
paper is not the best one available: width 128 at depth 4 reaches $9.7\pm0.5$
px\bp{} at $n{=}4$ against $18.2\pm4.2$ for our width 64, with four times as many
parameters, and we report the headline results at width 64 only for continuity
with the submitted version. Second, collocation is under-resolved at high
dimension: at $n{=}32$ raising $n_{\mathrm{col}}$ from $2000$ to $4000$ halves
the solve-error ($22.1$ to $11.0$ px\bp{}), which is the quadrature term of the
error estimate at work. Third, and this is the answer to where the variation
across $n$ comes from, the held-out residual is flat across $n$ to within twelve
per cent while the solve-error moves by up to fifty-seven per cent, with the same
rise to $n{=}16$ and fall at $n{=}32$ as the eight-seed sweep of
Table~\ref{tab:nconv}. The network solves the PDE equally well at every
dimension, so whatever variation there is in solve-error lies in the factor that
converts residual into solution error, the stability constant $C(n)$ of the
lifted operator \citep{mishra2023estimates,deryck2024error}---a property of the
Markovian lift, not of the optimiser. With two seeds per cell here and no
significant trend in the eight-seed sweep, we offer this as where the variation
lives, not as evidence that the error grows with $n$. The minutes-per-run column
of Table~\ref{tab:hpgrid} was recorded while other jobs shared the machine and
is indicative only; Table~\ref{tab:nconv} gives the controlled cost.
\begin{table}[!htbp]
\centering
\caption{Architecture grid at $n{=}4$, two seeds, $10^4$ iterations. Solve-error in px\bp{} against the 1M-path lifted MC; held-out PDE residual (RMS); minutes per run on CPU.}
\label{tab:hpgrid}
\small
\begin{tabular}{@{}r r r r r r@{}}
\toprule
width & depth & params & solve-error & residual & min \\
\midrule
\revrows{hparams_grid}
\bottomrule
\end{tabular}
\end{table}
\begin{table}[!htbp]
\centering
\caption{Held-out PDE residual across $n$ (width 64, depth 4, $n_{\mathrm{col}}{=}2000$, two seeds), with the solve-error and both quantities relative to $n{=}4$.}
\label{tab:hpnsweep}
\small
\begin{tabular}{@{}r r r r r@{}}
\toprule
$n$ & solve-error & residual & solve/solve$_4$ & resid./resid.$_4$ \\
\midrule
\revrows{hparams_nsweep}
\bottomrule
\end{tabular}
\end{table}
\begin{table}[!htbp]
\centering
\caption{Collocation ablation at $n{=}32$ (width 64, depth 4, two seeds).}
\label{tab:hpcolloc}
\small
\begin{tabular}{@{}r r r r@{}}
\toprule
$n_{\mathrm{col}}$ & solve-error & residual & min \\
\midrule
\revrows{hparams_colloc}
\bottomrule
\end{tabular}
\end{table}
\fi
\IfFileExists{../figures/fig_hparams.png}{\revtabletrue}{\revtablefalse}\ifrevtable
\begin{figure}[!htbp]\centering
\includegraphics[width=0.70\linewidth]{fig_hparams.png}
\caption{Solve-error (left axis) and held-out PDE residual (right axis) across the lift dimension $n$, width 64 / depth 4, two seeds.}
\label{fig:hparams}
\end{figure}
\fi

\subsection{Error decomposition and a design rule for $n$}
\label{sec:decomp_results}
The value of the decomposition \eqref{eq:split} is that it tells a user which
lever to pull, and the answer depends on the calibration. We have seen it both
ways in this work. At the parameters of the submitted version ($H{=}0.09$,
$\nu{=}0.30$), a single ``network versus true model'' error of $121$ vol\bp{} at
$n{=}4$ looked like a failure of the network; split by \eqref{eq:split} it was
almost entirely lift---$133$ vol\bp{} of lift-error against $37$ of
solve-error---and the lever was more factors. At the corrected calibration the
split reverses (implied-vol RMS, seed-averaged):
\begin{center}\small
\begin{tabular}{@{}r r r r@{}}
\toprule
$n$ & solve-error & lift-error & total (vol\bp) \\
\midrule
4 & $112$ & $76$ & $80$ \\
8 & $152$ & $33$ & $130$ \\
16 & $176$ & $9$ & $177$ \\
32 & $136$ & $26$ & $158$ \\
\bottomrule
\end{tabular}
\end{center}
The solve-error exceeds the lift-error at every $n$, so here the network is the
binding constraint and more factors buy nothing: the lever is the network
itself, and Section~\ref{sec:hparams} shows it is a real one, since a wider
network or twice the collocation roughly halves the solve-error. The two error
components also partly cancel at small $n$, which is why the total at $n{=}4$ is
smaller than the solve-error alone; a single end-to-end figure would have hidden
both the sign and the source of the error.

This gives a design rule that holds in both regimes: \emph{add lift factors only
while the lift-error exceeds the solve-error; beyond that, improve the network}.
The crossover is calibration-dependent. At the submitted parameters it fell near
$n{=}8$--$10$; at the corrected fit, whose smoother kernel ($H{=}0.17$) is easier
to approximate and whose stiffer factor dynamics are harder to solve, the
lift-error is below the solve-error already at $n{=}4$. A robustness sweep run
before this revision, at the submitted parameters and against a $400$-step
reference, found the crossover shifting with maturity in the expected direction:
the lift-error at $n{=}4$ grew from $8.6$ to $31.8$ px\bp{} as $T$ rose from
$0.05$ to $0.5$, so longer-maturity calibration needs more factors.

The network's error is not symmetric in strike. At every $n$ it over-prices at
and above the money, by $+138$ to $+221$ vol\bp{} ATM and $+130$ to $+203$ in the
calls, while the put wing is smaller and at $n{=}32$ near zero ($-15$). This is
a watch item for anyone deploying the operator inside a calibration loop, since
a one-signed error biases the recovered skew; we examine its mechanism next.
\IfFileExists{tables/n32bias.tex}{\revtabletrue}{\revtablefalse}\ifrevtable

\paragraph{The $n{=}32$ bias, taken apart}
Factor-tail coverage is symmetric by construction: the sampler widens both signs of the
factor deviation, and the asymmetric option is off by default, so asymmetry is not the
mechanism. Two others are separable by ablation: a noisy forward-variance anchor (the
baseline is anchored on a pre-flight Monte Carlo of $4\,000$ paths, and a level error in
the anchor is a level error in the price, which is what an ATM/call bias looks like), and
under-coverage of the factor tails in 32 dimensions. Table~\ref{tab:n32} re-estimates the
anchor with $200\,000$ paths and widens the tails, separately and together, against the
same reference of $10^6$ antithetic paths at $1\,600$ time steps, whose own
discretisation bias at this $n$ is $1.3$ px\bp{} (Table~\ref{tab:mcq}).
The ablation separates the two cleanly. A precise anchor lowers the solve-error
from $16.8$ to $14.3$ px\bp{} and the ATM bias by thirty vol\bp{}, so a noisy
anchor is part of the mechanism. Widening the tails does the opposite: it raises
the solve-error by $86$ per cent, roughly doubles the ATM bias and raises the
call bias by seventy per cent, alone or combined with the precise anchor. Wider tails spend collocation points
in regions of factor space the solution rarely visits, at the expense of the
bulk where the price is determined. The residual bias that survives a precise
anchor is the growth with $n$ that Section~\ref{sec:hparams} attributes to the
stability constant of the lifted PDE. All figures in this table are scored
against a $1\,600$-step reference; at the $400$ steps originally used, the
reference itself priced high by $3$--$4$ px\bp{} at this $n$ and masked part of
the network's over-pricing.
\begin{table}[!htbp]
\centering
\caption{Ablation of the $n{=}32$ directional bias, three seeds. Solve-error in px\bp; signed implied-vol error by strike bucket in vol\bp{} (network minus lifted MC).}
\label{tab:n32}
\footnotesize
\begin{tabular}{@{}l r r r r@{}}
\toprule
Configuration & solve-error & puts & ATM & calls \\
\midrule
\revrows{n32bias}
\bottomrule
\end{tabular}
\end{table}
\fi
\IfFileExists{../figures/fig_signed.png}{\revtabletrue}{\revtablefalse}\ifrevtable
\begin{figure}[!htbp]\centering
\includegraphics[width=0.75\linewidth]{fig_signed.png}
\caption{Signed solve-error by strike (network minus lifted MC, implied-vol basis points), seed-averaged, at $n{=}4,8,16,32$. The put wing stays near zero while the ATM and call errors grow with $n$.}
\label{fig:signed}
\end{figure}
\fi

\subsection{Solving the inverse problem: the correlation is pinned, the roughness is not}
\label{sec:btc}
Calibrated to the full Deribit BTC surface with a full-budget optimiser
(differential evolution, 60 generations, all 13 maturities and all 104 points,
followed by a bounded Nelder--Mead polish against the true Fourier forward
model), rough Heston returns
\begin{equation*}
  H{=}0.172,\quad \rho{=}{-}0.426,\quad V_0{=}0.105,\quad
  \theta{=}0.245,\quad \lambda{=}4.94,\quad \nu{=}1.37,
\end{equation*}
at a vega-weighted implied-vol RMSE of $113$ vol\bp{} and an unweighted RMSE of
$416$ vol\bp{} (Figure~\ref{fig:calib}). The unweighted figure is dominated by
one part of the surface. On the thirty-two points with maturities beyond two
months the model fits to $52$ vol\bp{}, comfortably inside a quoted Bitcoin
option spread; on the forty-eight points under two weeks it misses by $608$.
Rough Heston describes the BTC term structure well and fails only on the
shortest-dated options, whose smiles are too steep for this model to follow.
That is a statement about the forward model's expressiveness, not about the AI
forward operator, and it is orthogonal to the solve-error studied above.

\paragraph{A correction to the submitted fit}
The submitted version reported $H{=}0.090$, $\rho{=}{-}0.792$, $\nu{=}0.300$ at
$496$ vol\bp{}. That fit was produced by an objective with a defect we found
during this revision, and we state it plainly because it changes the headline
parameters. When the Fourier price at a point fell at or below intrinsic value,
the implied-volatility inversion returned no value and the objective silently
dropped the point and renormalised over the rest. A parameter vector at which
badly fitted points \emph{failed to price} therefore scored better than one that
priced them, and the optimiser could lower the loss by breaking the pricer. The
submitted $H{=}0.090$ sits in a window of width $2\times10^{-4}$ in $H$ created in
exactly that way. A second defect compounded it: the Lewis integrand decays like
$\exp(-VTu^2/2)$, so its truncation frequency must scale with $1/\sqrt{VT}$, and a
fixed truncation mis-priced the shortest maturities by up to $0.33$ in implied
volatility. Both are corrected. Every point now counts in every evaluation, a
point priced at or below intrinsic enters at the floor of the inversion bracket
so that its full error is charged, and a genuine pricer failure makes the whole
parameter vector infeasible rather than lighter. Scored under the corrected
objective the submitted vector reaches $204$ vol\bp{} weighted against $113$ for
the fit above.

The correlation and the vol-of-vol moved substantially; the roughness did not
move in the way that matters. Section~\ref{sec:roughness_robustness} shows that
this surface pins $\rho$ to within a single grid point but leaves $H$ free over
$[0.04,\,0.30]$ at five per cent of the best fit, so the point value $0.172$
above is one draw from that interval and not a measurement.

\begin{figure}[!htbp]\centering
\includegraphics[width=0.98\linewidth]{fig2_calibration_fit.png}
\caption{Rough Heston calibrated to the Deribit BTC surface (all 104 points, 13
maturities), solving the inverse problem with the true Fourier forward model.
(a)~Market versus model implied volatility at every point (dashed: identity;
colour: maturity); the outliers are the shortest maturities. (b)~Per-maturity
smiles for four representative maturities: market (markers) versus model
(dashed). Beyond two months the model fits to $52$ vol\bp{}; the residual is
concentrated in options under two weeks and reflects a model-expressiveness
limitation, not a forward-operator artefact.}
\label{fig:calib}
\end{figure}

\subsection{Is the recovered roughness real, and is it stable?}
\label{sec:roughness_robustness}
Whether Bitcoin volatility is rough rests, so far, on one option surface read
through a six-parameter inverse problem. Two things could undermine the claim.
The inverse problem might not identify $H$ sharply, in which case any single
fitted value would be one point on a flat ridge rather than a measurement; and
even if it does, a single snapshot says nothing about whether the roughness is a
stable feature of the market or an artefact of one day's regime. We test both,
the second using an estimator that never touches an option price.

\paragraph{Whether the surface identifies $H$}
It does not. The first evidence is simply that three full-budget fits of the
same surface under the same corrected objective, differing only in the search
box, disagree about $H$ by a factor of nine while agreeing about everything a
reader would check:
\begin{center}\small
\begin{tabular}{@{}l r r r r r@{}}
\toprule
search box & $H$ & $\nu$ & $\rho$ & weighted & unweighted \\
\midrule
original, unbounded polish & $0.019$ & $0.75$ & $-0.419$ & $114.9$ & $453.7$ \\
widened                    & $0.096$ & $0.98$ & $-0.421$ & $113.7$ & $430.2$ \\
widened further            & $0.172$ & $1.37$ & $-0.426$ & $113.2$ & $416.4$ \\
\bottomrule
\end{tabular}
\end{center}
The weighted objective differs by $1.7$ vol\bp{} across the three, and $\rho$
by $0.007$. The loss surface has a long, nearly flat ridge in $H$ along which
$\nu$ compensates: a rougher kernel with a smaller vol-of-vol prices this surface
about as well as a smoother kernel with a larger one.
A profile likelihood makes this precise (Table~\ref{tab:profile},
Figure~\ref{fig:profile}). Fixing $H$ on a twelve-point grid and re-optimising
the other five parameters from both ends of the ridge, the weighted objective
moves by only $1.3$ vol\bp{} as $H$ runs from $0.04$ to $0.26$, and every value
of $H$ in $[0.04,\,0.30]$ fits within five per cent of the best. The mechanism
is visible in the re-optimised parameters: $\nu$ climbs from $0.80$ to $1.94$
and $\lambda$ from $2.3$ to $7.6$ along the ridge while $\rho$ barely moves.
Profiling $\rho$ in the same way gives the opposite picture: the objective rises
steeply on either side of $\rho=-0.45$ ($146.5$ vol\bp{} at $-0.60$, $138.4$ at
$-0.30$, against $115.0$ at $-0.45$), so the five-per-cent set is a single grid
point. \emph{The surface identifies the correlation sharply and the roughness
barely at all.} Because the whole $H$ interval lies far below $\tfrac12$, the
qualitative claim that Bitcoin volatility is rough survives; the point estimate
does not, and we do not report one. The profile's minimum, $113.2$ vol\bp{} at
$H=0.18$, matches the global fit; because each point is a local re-optimisation,
the reported values are upper bounds and the interval is, if anything, slightly
too narrow. An earlier pass of this profile confined $\lambda$ to the box
$[0.05,3]$, where it sat on the upper bound at ten of twelve grid points; the
figures above use a box wide enough to contain it, which widened the $H$
interval from $[0.04,0.18]$ to $[0.04,0.30]$. One bound still binds, and it
binds at the edge that matters: from $H{=}0.30$ upward the re-optimised $\nu$
sits on its cap of $2$, so the profile there is itself box-limited and the upper
end of the interval is a lower bound on the true one. The surface may support
still smoother kernels than we report.
We therefore do three things. We report the option-implied $H$ as an interval
rather than a point. We run every PINN experiment in this paper at the best
fit, $H=0.172$ with $\nu=1.37$ and $\lambda=4.94$, and we note that this is a
deliberately demanding choice: the fast mean-reversion and large vol-of-vol make
the lifted factor system markedly stiffer than at the parameters of the
submitted version, and every solve-error in Section~\ref{sec:nconv} is larger
for it. And we lean on the estimate below, which never passes through this
inverse problem at all.

\IfFileExists{tables/profile.tex}{\revtabletrue}{\revtablefalse}\ifrevtable
\begin{table}[!htbp]
\centering
\caption{Profile likelihood in $H$ on the headline surface. $H$ is fixed at each grid value and the other five parameters are re-optimised by Nelder--Mead from both ends of the ridge, keeping the better. Weighted and unweighted implied-vol RMSE in vol\bp; the last three columns show how $\nu$, $\rho$ and $\lambda$ compensate along the ridge.}
\label{tab:profile}
\small
\begin{tabular}{@{}r r r r r r@{}}
\toprule
$H$ & weighted & unweighted & $\nu$ & $\rho$ & $\lambda$ \\
\midrule
\revrows{profile}
\bottomrule
\end{tabular}
\end{table}
\fi
\IfFileExists{../figures/fig_profile.png}{\revtabletrue}{\revtablefalse}\ifrevtable
\begin{figure}[!htbp]\centering
\includegraphics[width=0.90\linewidth]{fig_profile.png}
\caption{Profile likelihood on the headline surface. Left: $H$ fixed on a grid, the other five parameters re-optimised from both ends of the ridge; right: the same for $\rho$. Vertical line: the canonical fit; dotted line: 5\% above the best profile value, the interval rule.}
\label{fig:profile}
\end{figure}
\fi
\IfFileExists{tables/risk.tex}{\revtabletrue}{\revtablefalse}\ifrevtable
\paragraph{What the ridge costs in risk terms}
Non-identification of $H$ matters only if the quantities a desk actually uses depend on where
along the ridge one sits. Table~\ref{tab:risk} prices three of them with the exact forward
model at three months, for every parameter vector the surface cannot tell apart: the
variance-swap fair strike (from the log-contract replication), the 25-delta risk reversal,
and a 90/110 collar.
The answer is reassuring, and it follows from the mechanism above. Across the nine
ridge vectors---$H$ from $0.04$ to $0.30$, all within five per cent in fit---the
variance-swap fair strike moves by $0.60$ vol points on a level of $43.4$, and
the risk reversal by $0.46$; the collar, which is sensitive to the curvature of
the smile between its two strikes, moves by $5.7$ px\bp{}. A seven-fold
ambiguity in $H$ is therefore close to harmless for these functionals, because
$\nu$ and $\lambda$ compensate and the functionals depend on the compensated
combination rather than on $H$ alone. The practical corollary is the one a desk
needs: what this surface cannot tell apart, these instruments do not care about.
The three hybrid vectors of Section~\ref{sec:parametric} tell a different story,
and a useful one. They are fitted to the $T{=}0.15$ slice alone, and they price a
risk reversal about half a vol point steeper and a collar seven basis points
dearer than any full-surface fit: the ambiguity that matters in practice is not
where on the ridge a calibration lands, but which part of the surface it was
asked to fit.
\begin{table}[!htbp]
\centering
\caption{Risk functionals along the calibration ridge, Fourier-priced at $T{=}0.25$, $S_0{=}1$, $r{=}0$. Variance swap as $\sqrt{K_{\mathrm{var}}}$ in annualised vol points; risk reversal as IV(25$\Delta$ call)$-$IV(25$\Delta$ put) in vol points; collar as long 90\% put, short 110\% call, in price basis points of spot.}
\label{tab:risk}
\small
\begin{tabular}{@{}l l r r r@{}}
\toprule
Vector & & Var.\ swap & $25\Delta$ RR & Collar \\
\midrule
\revrows{risk}
\bottomrule
\end{tabular}
\end{table}
\fi

\paragraph{An estimate that does not use the option surface}
Following Gatheral, Jaisson and Rosenbaum~\citep{gatheral2018volatility}, we estimate $H$ directly from realised
volatility. From five-minute log returns of the Deribit BTC perpetual we build
daily realised variance, set $\sigma_t=\sqrt{RV_t}$, and form the $q$-th absolute
moment of log-volatility increments at lag $\Delta$,
$m(q,\Delta)=\mathbb{E}\bigl|\log\sigma_{t+\Delta}-\log\sigma_t\bigr|^{q}$. For a
fractional process $m(q,\Delta)\sim\Delta^{qH}$, so regressing $\log m$ on
$\log\Delta$ for each $q$ gives slopes $\zeta(q)$, and regressing $\zeta$ on $q$
through the origin gives $H$. Over $2\,921$ days (August 2018 to August 2026,
$841$k bars) this returns
\[
  H_{\mathbb{P}} = 0.111, \qquad \text{95\% CI } [0.102,\,0.150],
\]
with the interval from a $60$-day moving-block bootstrap. This is an independent
corroboration in the strongest available sense: it uses no option data, no
Markovian lift, no Fourier pricer and no optimiser, and it lands inside the
option-implied interval $[0.04,\,0.30]$. Where the option surface cannot choose a
point on its ridge, this estimate can, and it chooses one well inside the rough
regime. It also agrees with Takaishi~\citep{takaishi2020rough}, who
reports rough volatility for Bitcoin over an earlier and disjoint sample.

We report the block-bootstrap interval rather than the regression's own standard
error deliberately. The analytic standard error propagates only the scatter of
the five $\zeta(q)$ about the fitted line; it ignores the sampling error inside
each moment and the strong persistence of log-volatility, and on our data it is
roughly an order of magnitude too small---it returns $2\times10^{-4}$ for a
window whose neighbouring window sits $0.02$ away. Any inference drawn from it
would be spurious.

\paragraph{Stability across the spot-ETF approval}
The referee's regime question has a natural event: the approval of US spot
Bitcoin ETFs on 11 January 2024. We split the realised-volatility series into
calendar 2023 (pre) and February--December 2024 (post), excluding January 2024
entirely as the event window rather than assigning it to either regime, and
re-estimate:
\[
  H_{\mathbb{P}}^{\text{pre}} = 0.064\ [0.018,\,0.098], \qquad
  H_{\mathbb{P}}^{\text{post}} = 0.062\ [0.022,\,0.079],
\]
a difference of $+0.002$ against a combined bootstrap standard error of $0.025$
($z=0.10$). The roughness of Bitcoin volatility is unchanged by the ETF
approval. One caveat belongs with these numbers: this estimator is biased low on
short windows, so the two one-year windows are comparable with each other but not
with the full-sample figure above, and we compare only the former.
\IfFileExists{tables/regime.tex}{\revtabletrue}{\revtablefalse}\ifrevtable

\paragraph{Option-implied, across the same split}
The same event study can be run on option surfaces, with the caveat of
Section~\ref{sec:data}: the historical panel is trade-based where the headline snapshot is
quote-based. From it we build 52 pre-ETF (calendar 2023) and 47 post-ETF (February--December
2024) weekly surfaces. Because a single surface does not pin $H$, each week is fitted by
Nelder--Mead from both ends of the ridge, never from the previous week, and
Table~\ref{tab:regime} compares three things across the split: $H$ within each basin, the
identified parameters $\rho$ and $\nu$, and which basin the week prefers.
The roughness does not move. From either warm start the mean weekly $H$ is
statistically indistinguishable before and after the approval (canonical start
$0.099$ against $0.107$, bootstrap interval on the difference $[-0.071,+0.059]$;
wide start $0.163$ against $0.132$, $[-0.059,+0.092]$), which agrees with the
physical-measure result above from an entirely independent route. Two things do
move: $\nu$ falls from $1.96$ to $1.63$ and the fit degrades from $124$ to $151$
vol\bp{}, both with bootstrap intervals that exclude zero, and $\rho$ shifts by
$-0.083$ with an interval that only just excludes it. We read the last with
care. The weekly $\rho$ sits near zero, well away from the $-0.45$ of the
headline snapshot, and that gap is the trade-versus-quote provenance difference
of Section~\ref{sec:data} rather than a property of the market: a comparison
\emph{across} weeks of the panel is sound, a comparison of its level with the
headline fit is not. The weekly $H$ itself varies widely from week to week and
between the two starts on the same week, which is the non-identification of the
headline surface reappearing at weekly resolution; the stability test is
meaningful only because each regime is summarised over roughly fifty weeks.
\begin{table}[!htbp]
\centering
\caption{Option-implied parameters before and after the spot-ETF approval, 52 and 47 weekly surfaces. Mean\,$\pm$\,SD per regime; difference pre$-$post with a 4-week block-bootstrap 95\% interval and Welch $p$.}
\label{tab:regime}
\footnotesize
\begin{tabular}{@{}l r r r r r@{}}
\toprule
Quantity & pre-ETF & post-ETF & diff. & 95\% CI & $p$ \\
\midrule
\revrows{regime}
\bottomrule
\end{tabular}
\end{table}
\fi
\IfFileExists{../figures/fig_regime.png}{\revtabletrue}{\revtablefalse}\ifrevtable
\begin{figure}[!htbp]\centering
\includegraphics[width=0.90\linewidth]{fig_regime.png}
\caption{Weekly option-implied $H$ from the trade-tape panel, fitted from both ends of the ridge (top), and which end each week prefers (bottom: loss from the wide-box start minus loss from the canonical start, negative preferring the smoother $H$). Shaded: January 2024, excluded as the ETF-approval window.}
\label{fig:regime}
\end{figure}
\fi

\subsection{A parametric AI forward operator inside the inverse problem}
\label{sec:parametric}
A single network cannot cover both the breadth of a parameter box and the
accuracy of a point-trained forward operator. Table~\ref{tab:calib} and
Figure~\ref{fig:hybrid} report the calibration experiment with three seeds.

The parametric network adds $(H,\nu,\rho,\lambda,\theta)$ as five additional
inputs, holding $V_0$ and the maturity fixed; we use $n{=}10$ (the crossover
rule, Section~\ref{sec:decomp_results}) and width 96, depth 5. We train it on two
boxes with identical architecture, budget and seeds: the original box of the
submitted version ($\nu\le0.6$, $\theta\le0.35$) and a widened one ($\nu\le1.0$,
$\theta\le0.45$) that contains the optimum of the $T{=}0.15$ slice. That turns
the question of whether the box was mis-specified into a controlled ablation.

Solve-error depends strongly on where in parameter space it is measured, and
the probe points have to be read with that in mind. At mild parameter vectors
the network is accurate---$3.6$ to $9.1$ px\bp{} at the weak-skew and
strong-skew probes---while at the rough, low-$H$ probe it reaches $32$--$38$
px\bp{}. At the calibrated vector itself the widened box gives $9.1\pm3.4$
px\bp{} against $33.8\pm2.5$ for the original; that vector has $\nu=1.37$, outside
\emph{both} training boxes, so both figures are extrapolations and the widened
box's advantage is that it extrapolates a shorter distance. The corners are
quantified in Figure~\ref{fig:heatmap}: the error concentrates where $H$ is small
and $\nu$ large, reaching $38$--$40$ px\bp{}, and it is an interaction rather
than the sum of two effects---adding a $\nu/H$ term to an additive regression
raises $R^2$ from $0.42$--$0.49$ to $0.67$--$0.68$, a partial $R^2$ of $0.19$ to
$0.25$. That is the generalisation limit of this AI forward operator, stated
explicitly rather than averaged away.
\IfFileExists{../figures/fig_heatmap.png}{\revtabletrue}{\revtablefalse}\ifrevtable
\begin{figure}[!htbp]\centering
\includegraphics[width=0.90\linewidth]{fig_heatmap.png}
\caption{Parametric solve-error over $(H,\nu)$ at the box centre in $(\rho,\lambda,\theta)$, $n{=}10$, seed-averaged, against a lifted Monte Carlo at every cell. Titles give the interaction coefficient $d$ of $\mathrm{err}\sim a+b/H+c\nu+d\,\nu/H$ and its partial $R^2$.}
\label{fig:heatmap}
\end{figure}
\fi

Used as a standalone inverse-problem solver the parametric operator is fast
(approximately $2$\,s, zero true-model calls) but poor: $161\pm15$ vol\bp{} under
the true model on the original box and $181\pm97$ on the widened one. The
submitted version attributed this to the slice optimum lying outside the
training box. The ablation does not support that explanation. The widened box
contains the optimum ($\nu=0.99$ against a bound of $1.0$), yet its standalone
fit is no better on average and far less reliable---the three seeds reach $316$,
$92$ and $134$ vol\bp{}. Containing the optimum is necessary but not
sufficient: the network fits its own imperfect forward map rather than the
market, and over a wider box that map is less accurate where it matters. We do
not recommend the standalone parametric operator as an inverse-problem solver,
on either box.

Used as a warm start it is effective, and on both boxes. Seeding a Nelder--Mead
polish against the true forward model with the network's parameter guess reaches
$23.8$ vol\bp{} from the original box and $23.6$ from the widened one, against
$22.7$ for the from-scratch fit---on a slice whose market implied vols span
roughly $34$ to $52$ per cent, so well inside a quoted spread. It does so in
$197$ and $205$ true-forward-model evaluations against $1\,736$, a saving of
$8.8\times$ ($8.7$--$8.9\times$ across seeds) and $8.5\times$ respectively, and in
$17$\,s against $142$\,s. The warm start is insensitive to which box the network
was trained on, which is the useful practical conclusion: the hybrid does not
require the training box to be chosen well. The value proposition is amortised: train the AI
forward operator once, then solve each new instance of the inverse problem
cheaply, with the exact forward model kept in the loop for a brief, decisive
finish. The amortisation has a concrete break-even. Each warm-started
recalibration saves about $1\,535$ true-model evaluations, or $125$\,s of
wall-clock; training the parametric network costs $87$ to $110$ minutes on the
widened and original boxes. The network pays for itself after $42$ to $53$
recalibrations of the same maturity slice: about two months of daily
recalibration, or three and a half to four and a half hours of recalibration
every five minutes. Below
that, the exact forward model from scratch is the cheaper tool, and we would
recommend it. The training times were measured while other jobs shared the
machine, which inflates them, so the break-even is if anything overstated.

\begin{table}[!htbp]
\centering
\caption{Solving the calibration inverse problem for the BTC $T{\approx}0.15$
slice: the true-forward-model loop from scratch, and the AI-only and hybrid
routes with the parametric network trained on the original box and on a widened
box that contains the slice optimum. Three network seeds each; fit quality is
judged by the true Fourier forward model over the same feasible region.}
\label{tab:calib}
\small
\begin{tabularx}{\linewidth}{@{}>{\raggedright\arraybackslash}X c c c@{}}
\toprule
Method & Fit RMSE (vol\bp) & Wall-clock (s) & True-model evals \\
\midrule
%% rows come from make_tables.py (results/param_boxes.json) once the w96/d5 nets on both
%% boxes are evaluated; the hand-typed w80/d4 single-box rows below are the fallback
\IfFileExists{tables/calib.tex}{\revtabletrue}{\revtablefalse}\ifrevtable
\revrows{calib}
\else
True-forward-model loop, from scratch & $23.7$       & $168$      & $2018$     \\
AI-only (single-shot)                 & $222\pm20$   & ${\sim}2$  & $0$        \\
Hybrid (AI warm start, true-model polish) & $29.4\pm7.1$ & $23$   & $237\pm46$ \\
\fi
\bottomrule
\end{tabularx}
\end{table}

\begin{figure}[!htbp]\centering
\includegraphics[width=0.82\linewidth]{fig3_hybrid.png}
\caption{Hybrid solution of the inverse problem. Fit quality (left) and
true-forward-model evaluation count (right) for the three methods. The hybrid
matches from-scratch quality at $8.5$--$8.8\times$ fewer true-model evaluations
whichever box the network was trained on; the AI-only route is fast but
inaccurate, and widening the box to contain the optimum makes it no better on
average and far less reliable across seeds.}
\label{fig:hybrid}
\end{figure}

\IfFileExists{tables/surrogate.tex}{\revtabletrue}{\revtablefalse}\ifrevtable
\subsection{Against the forward operators practitioners use}
\label{sec:operators}
A reader who wants a fast forward operator for rough Heston would not reach for a PINN first.
They would price with the Fourier transform directly, or run a short lifted Monte Carlo, or
generate labels with an exact pricer and fit a network to them in the manner of
Bayer et al.~\citep{bayer2019deep} and Horvath et al.~\citep{horvath2021deep}---the method class of the calibration
networks of Liu et al.~\citep{liu2019cann}. Table~\ref{tab:operators} places our operator next to all
three under one protocol. The supervised surrogate is built to match the parametric network
exactly: the same five inputs, the same fixed 21-point log-moneyness grid at $T{=}0.15$, the
same widened box, Sobol-sampled and split by parameter vector so that no strike of a test
vector is ever seen in training.
We state the comparison plainly, because it goes against us on the axes most
readers will check first. The supervised surrogate, trained on $64\,353$
Fourier-labelled parameter vectors, reaches $5.7\pm0.1$ vol\bp{} against the
transform method, and $9.3\pm1.0$ in the rough, high-$\nu$ corner where our
parametric network degrades most. It is more accurate than the PINN against
Fourier, and once trained it is as fast. On one-shot cost both the Fourier
pricer and a short lifted Monte Carlo beat a PINN training run by three to four
orders of magnitude. What the surrogate cannot do is exist without an exact
pricer: every one of its labels came from the fractional Riccati solver, and in
a model with no transform method it cannot be built at all. The PINN needs only
the PDE. It also yields the solve/lift decomposition of
Section~\ref{sec:decomp}, which a network fitted to Fourier labels inherits the
lift error from silently and cannot separate. Those two properties---no exact
pricer offline, and an error budget split into what training controls and what
the lift dimension controls---are the case for this operator, and we make no
claim beyond them.
\begin{table}[!htbp]
\centering
\caption{Forward operators for the lifted rough-Heston smile at $T{=}0.15$.
Accuracies are not comparable across rows, and the column says against what each
is measured: the surrogate is trained on Fourier labels and scored against them,
in implied-vol basis points, while the PINNs are scored in price basis points
against an independent lifted Monte Carlo at the same $n$ and carry the
lift-error against Fourier on top. The parametric range runs from mild
parameter vectors to the rough, high-$\nu$ corner. Both PINNs and the surrogate
are differentiable in their inputs; the Fourier pricer and the Monte Carlo are
not. Costs for the PINNs are CPU wall-clock; the parametric figure was measured
on a shared machine and is indicative.}
\label{tab:operators}
\footnotesize
\begin{tabularx}{\linewidth}{@{}>{\raggedright\arraybackslash}X >{\raggedright\arraybackslash}p{2.3cm} >{\centering\arraybackslash}p{1.5cm} >{\raggedright\arraybackslash}p{2.7cm} >{\centering\arraybackslash}p{1.5cm}@{}}
\toprule
Operator & Offline cost & Exact pricer offline? & Accuracy (against) & Solve/lift split? \\
\midrule
Fourier (fractional Riccati) & none & is one & exact by construction & no \\
Lifted MC, $10^6$ paths & none & no & reference; SE $0.7$--$0.8$ px\bp & no \\
Single-point PINN, $n{=}4$ & ${\approx}20$ min per vector & no & $15.2$ px\bp{} (lifted MC) & yes \\
Parametric PINN, $n{=}10$ & $80$--$125$ min, once & no & $3.6$--$38$ px\bp{} (lifted MC) & yes \\
\revrows{surrogate}
\bottomrule
\end{tabularx}
\end{table}
\fi

% ================================================================
\section{Discussion}
\label{sec:disc}

\paragraph{Convergence of the AI forward operator}
The finding we care about most is that the network's own error in solving the
lifted forward PDE does not measurably degrade as the dimension grows, while
grid methods become impractical as an inverse-problem engine. The history of that
claim is worth recording, because it moved twice. An early single-seed run made
the error look flat; three seeds suggested a doubling from $n{=}4$ to $32$; eight
seeds, a corrected reference and a formal test give a log-log slope of
$0.09\pm0.08$ and no rejection of flatness ($p=0.14$). We state the result in the
form the eight-seed data support---no detectable growth, and growth that is in
any case significantly sub-linear---and we note that it holds at the stiffer
corrected calibration, where the absolute error is also larger. The grid
comparison uses the naive $50^{n+1}$ count; sparse-grid, ADI and tensor-train
methods go further, so the defensible claim is \emph{impractical, not
impossible}. Even so, the contrast---an error and a cost that do not grow
measurably with $n$, against a grid that cannot be stored---is a concrete
instance of the mesh-free advantage as a component of an inverse-problem
pipeline.

\paragraph{Stability: four design choices are load-bearing}
We ablated each stabiliser explicitly. Without the hard-constrained ansatz, the
soft IC penalty oscillates against the residual and the forward operator fails
to converge. Without the integrated-variance anchor, a uniform $120$ vol\bp{}
level bias remains and propagates directly into every downstream calibration.
Without MC-informed factor sampling, the smile flattens and RMSE roughly
doubles. Without self-adaptive weighting or the $\tau$-curriculum, seed variance
is markedly higher (from $339\pm62$ to $310\pm28$ vol\bp{} in an early stability
comparison). Each of these choices carries over to other AI forward
operators for high-dimensional inverse problems.

\paragraph{A lesson about validating AI forward operators}
The error decomposition revealed something methodologically important. When a
learned surrogate replaces an exact forward solver inside an inverse problem, it
is dangerously easy to validate it against the wrong reference. Our network
solves the finite-$n$ lifted PDE, whose exact solution is the lifted Monte
Carlo, not the $n\to\infty$ true model; comparing straight to the true model
conflates network error with model error. We found the two carry opposite signs
and that the historical ``network-versus-true-model'' gap was lift-dominated,
not a failure of the AI method. \emph{Validate an AI forward operator against a
same-fidelity reference, report the two errors separately, and let their
crossover set the approximation order}: none of this is special to rough
volatility, and we urge it on anyone placing a learned forward operator inside
an inverse problem built on an approximated governing equation.

\paragraph{Generalisation and identifiability inside the inverse problem}
A single network spread over five parameters is markedly less accurate than a
point-trained one, worst in the rough, high-vol-of-vol corners where the PDE is
stiffest---an explicit, quantified generalisation limit rather than an averaged
headline number. Used alone as an inverse-problem solver it is additionally
vulnerable to a training-box identifiability failure: if the true optimum lies
outside the box the parametric operator was trained on, no amount of gradient
descent inside that box recovers it. This is not a limitation of the inverse
problem itself, which remains well posed once regularised by the Fourier forward
model; it is a limitation of trusting the AI forward operator unsupervised. The
warm-start hybrid sidesteps exactly this failure mode by keeping the true
forward model as the final arbiter.

\paragraph{Limitations and watch items}
Our evidence rests on a single market snapshot (one Deribit BTC surface), and
the parametric and calibration studies fix the maturity to a single slice.
Timings are single-machine CPU figures at moderate budgets. The mild directional
solve-bias at large $n$ (ATM and call over-pricing growing to ${\sim}100$
vol\bp{} at $n{=}32$) is a watch item for high-$n$ deployments inside a
calibration loop. The sub-linear trend need not continue arbitrarily far.

\paragraph{Future work}
Natural extensions include: (i)~adding the maturity as a network input so a
single forward operator serves the whole surface's inverse problem at once;
(ii)~scaling width, depth, and GPU training to shrink the parametric corner
penalty and widen the training box, directly addressing the identifiability
failure mode above; (iii)~extending to other rough models (rough Bergomi) and to
American and path-dependent payoffs, where the inverse problem is harder still;
(iv)~investigating whether residual reweighting or a correction term resolves
the high-$n$ directional bias; (v)~a formal stability/identifiability analysis of
the calibration inverse problem itself as $n\to\infty$, complementing the
purely empirical treatment given here.

% ================================================================
\section{Conclusion}
\label{sec:conc}

We have characterised a physics-informed neural network as a forward operator
for the rough-volatility calibration inverse problem, on the
$(n{+}1)$-dimensional Markovian-lifted PDE that makes the forward problem
well-posed in the first place. Across an eight-fold rise in the lifted
dimension the network's solve-error shows no statistically detectable growth
and its training cost per iteration is flat, while naive finite-difference
forward solvers become infeasible by $n{=}8$; we identify the four training
choices that make this behaviour hold, and show that the held-out PDE residual
is equally flat, so what variation remains lies in the stability constant of the
lifted operator rather than in the optimiser. Separating solve-error from
lift-error is essential to correctly attribute inverse-problem error: it told us
to add factors at one calibration and to improve the network at another, and it
yields a design rule for the factor count that holds in both. Applied to live
Bitcoin options, the surface pins the correlation but bounds the Hurst exponent
only to $[0.04,0.30]$---rough, but not a point---and that ambiguity moves the
variance-swap strike and the risk reversal by well under a vol point. A
parametric extension of the forward operator, while not accurate enough to solve
the inverse problem independently, is an effective warm start that solves it at
$8.8$ times fewer expensive forward-model evaluations. The methodological message is that an AI forward operator earns
its place inside an inverse-problem pipeline only once its own error is
separated from the model error it is approximating---and that, done this way,
it degrades gracefully where grid methods cannot follow.

% ================================================================
\section*{Declaration of competing interest}

The authors declare that they have no known competing financial interests or
personal relationships that could have appeared to influence the work reported in
this paper.

\section*{Funding}

This research received no specific grant from any funding agency in the public,
commercial, or not-for-profit sectors.

\section*{Data availability statement}

The BTC option data were retrieved from the Deribit public REST API
(no API key required) at \texttt{https://www.deribit.com/api/v2}. The raw data
snapshot used in the experiments, the calibrated parameter files, the two
robustness datasets of Section~\ref{sec:data} in derived form (weekly implied-vol
surfaces and daily realised variance), every result file the tables are generated
from, and the complete PyTorch implementation are available at
\url{https://github.com/aviral9876/physics-informed-nn-rough-vol}.
%% DOI: the tagged release v1.0-cnsns is to be archived on Zenodo; insert the DOI
%% sentence here once minted (needs the corresponding author's account).

\section*{Declaration of generative AI and AI-assisted technologies in the
writing process}

During the preparation of this manuscript the authors used Claude (Anthropic)
to assist with prose drafting and structural organisation. After using this
tool, the authors reviewed and edited all content as needed and take full
responsibility for the accuracy and integrity of the published work. The
mathematical derivations, experimental results, code, and all scientific claims
are the authors' own.

% ================================================================
\appendix

\section{Derivation of the lifted pricing PDE}
\label{app:derivation}

We derive \eqref{eq:pde} from the lifted dynamics \eqref{eq:lift}, since a
referee correctly observed that the correspondence is not self-evident --- in
particular the presence of a \emph{full} $(i,j)$ double sum in the factor block,
where a reader expecting independent Ornstein--Uhlenbeck factors would predict a
diagonal one.

Collect the state as $Z_t=(X_t,U^1_t,\ldots,U^n_t)$ with $X_t=\log S_t$. Under
the risk-neutral measure the lifted system is
\begin{align}
  dX_t   &= \left(r-\tfrac12 V_t\right)dt + \sqrt{V^+_t}\,dB_t, \\
  dU^i_t &= \left[-x_i U^i_t + \lambda(\theta-V_t)\right]dt
            + \nu\sqrt{V^+_t}\,dW^V_t, \qquad i=1,\ldots,n,
\end{align}
with $V_t=V_0+\sum_i c_i U^i_t$ and $d\langle B,W^V\rangle_t=\rho\,dt$. This is
a $(n{+}1)$-dimensional It\^o diffusion, hence Markovian --- which is the entire
purpose of the lift, the original Volterra variance process being neither.

The instantaneous covariation matrix of $Z$ follows from the fact that the
system is driven by only \emph{two} Brownian motions, $B$ and $W^V$, whatever
the value of $n$:
\begin{equation}
  \begin{split}
  \frac{d\langle X,X\rangle_t}{dt}&=V^+_t,
  \qquad
  \frac{d\langle X,U^i\rangle_t}{dt}=\rho\nu V^+_t,\\
  \frac{d\langle U^i,U^j\rangle_t}{dt}&=\nu^2 V^+_t
  \quad\text{for \emph{all} }i,j.
  \end{split}
\end{equation}
The last identity is the crux. Every factor is driven by the same increment
$dW^V_t$ and differs only through its mean-reversion speed $x_i$, so the factor
block of the covariation matrix is $\nu^2V^+_t\,\mathbf{1}\mathbf{1}^{\!\top}$
--- rank one, with every entry equal, rather than diagonal. The factors are
perfectly instantaneously correlated; they decorrelate only through their
differing decay rates over finite time. Consequently the second-order factor
term in the generator carries the unrestricted double sum
$\sum_{i,j=1}^{n}P_{u_iu_j}$ that appears in \eqref{eq:pde}, and not
$\sum_{i}P_{u_iu_i}$.

The infinitesimal generator of $Z$ is therefore
\begin{align}
  \mathcal{A}P = {}& \left(r-\tfrac12 V\right)P_x
      + \sum_{i=1}^{n}\left[-x_i u_i+\lambda(\theta-V)\right]P_{u_i}
      \nonumber\\
    & + \tfrac12 V P_{xx}
      + \tfrac12\nu^2 V\sum_{i,j=1}^{n}P_{u_iu_j}
      + \rho\nu V\sum_{i=1}^{n}P_{xu_i},
\end{align}
where we write $V$ for $V^+$ to lighten notation, the truncation being active
only on a set the scheme controls numerically. For a European claim with payoff
$g(X_T)$, the discounted price $P(t,x,u)=\mathbb{E}\!\left[e^{-r(T-t)}g(X_T)\mid
Z_t=(x,u)\right]$ satisfies the Feynman--Kac equation
$P_t+\mathcal{A}P-rP=0$ with $P(T,x,u)=g(x)$, which on substituting
$\mathcal{A}$ and $g(x)=\max(e^{x}-K,0)$ is exactly \eqref{eq:pde}.

Three remarks bear on the numerics. First, the drift source
$\lambda(\theta-V)$ and the diffusion coefficient $\nu\sqrt{V}$ are
\emph{shared} across factors; only $-x_iu_i$ distinguishes them, which is why
the factors can be advanced with a single common Brownian draw. Second, the
$(x_i)$ span roughly ten orders of magnitude at small $H$, making the factor
block stiff and forcing the exponential-integrator treatment of
Section~\ref{sec:eval}. Third, the rank-one structure of the factor covariation
is what couples $\rho$ and $\nu$ across the whole factor block, and is the
reason the low-rank tensor structure that a tensor-train solver would need is
eroded rather than preserved as $n$ grows.

% ================================================================
\bibliographystyle{elsarticle-num}
\bibliography{refs}

\end{document}
